\section{Algorithms and Code}
\label{sec:algorithms}
Algorithms and implementation details should be presented at the level
needed for the reader to understand, reproduce, or evaluate the work.
Use pseudocode when the logic of an algorithm is the important part, source
code listings when a specific implementation detail matters, and equations
when the underlying computation or relationship needs to be stated
precisely. Avoid reproducing implementation details that do not contribute
to the argument; the complete source code can instead be provided in an
appendix or repository.

\subsection{Pseudocode}
When the \emph{logic} of a method matters more than a specific programming
language, present it as pseudocode with the \texttt{algorithm2e} package,
as in \Cref{alg:kmeans}. Pseudocode should be language-agnostic, use
meaningful variable names, and be short enough to read in under a minute.
If it needs more than about 20 lines, summarise the outer loop in prose and
put the full version in an appendix.

\begin{algorithm}[htbp]
  \DontPrintSemicolon
  \KwIn{Dataset $X = \{x_1,\ldots,x_n\}$, number of clusters $k$}
  \KwOut{Cluster assignments $c_1,\ldots,c_n$ and centroids $\mu_1,\ldots,\mu_k$}

  Initialise $k$ centroids $\mu_1,\ldots,\mu_k$ randomly\;
  
  \Repeat{convergence}{
    \ForEach{$x_i \in X$}{
      $c_i \gets \arg\min_{j \in \{1,\ldots,k\}}
      \lVert x_i - \mu_j \rVert^2$\;
    }

    \ForEach{$j \in \{1,\ldots,k\}$}{
      $\mu_j \gets \frac{1}{|\{i : c_i=j\}|}
      \sum_{i:c_i=j} x_i$\;
    }
  }

  \Return{$c_1,\ldots,c_n,\mu_1,\ldots,\mu_k$}\;

  \caption{K-means clustering}
  \label{alg:kmeans}
\end{algorithm}

\subsection{Source code listings}
When the exact code matters (e.g.\ when explaining a preprocessing step or
a model implementation detail), use the \texttt{listings} package rather
than a screenshot, so the code is searchable, copy-pasteable, and
consistently styled. \Cref{lst:kmeans} shows a short Python implementation
using scikit-learn.

\begin{lstlisting}[language=Python,
  caption={Clustering data points with K-means.},
  label={lst:kmeans}]
from sklearn.cluster import KMeans

model = KMeans(n_clusters=3, random_state=42)
labels = model.fit_predict(X)
\end{lstlisting}

\begin{tipbox}
  Only include code that illustrates a specific point you are making in
  the text. A report is not the place to reproduce your entire codebase
  --- link to a repository or put long listings in an appendix.
\end{tipbox}

\begin{pitfallbox}
  Do not paste an indented code block as plain text without syntax
  highlighting or numbering; it becomes very hard to reference (``see line
  14'') and hard to read. Use \texttt{listings} or \texttt{minted}
  consistently throughout the report.
\end{pitfallbox}

\subsection{Equations}
Number equations you refer to later, and treat the equation as part of
the sentence grammatically. For example, the K-means objective can be
defined as

\begin{equation}
  \label{eq:kmeans}
  J =
  \sum_{i=1}^{n}
  \left\lVert x_i - \mu_{c_i} \right\rVert^2,
\end{equation}

where $x_i$ is the $i$th data point, $c_i$ is its assigned cluster, and
$\mu_{c_i}$ is the centroid of that cluster. \Cref{eq:kmeans} expresses
the objective that K-means minimises: the total squared distance between
each data point and the centroid of its assigned cluster. The equation
therefore gives the reader a precise mathematical description of the
algorithm presented in \Cref{alg:kmeans}.

